\documentclass{article} %TODO refs

\usepackage[utf8]{inputenc}
\usepackage{amsmath}
\usepackage{graphicx}

\title{Question justifications}
\author{Maham S}
\date{\today}

\begin{document}
\maketitle
\newpage
\tableofcontents
\newpage
\section{Pre-session Questionnaire}
\subsection{\textbf{Participant Experience Information}}
The first few questions of this survey are used to confirm the identity of the participant being interviewed. \\
The next few questions examine the participant's experiences. '5 years' is a consertive amount of time that was chosen, because technology is ever-changing, 
and is projected to change in the near future (https://www.hunimed.eu/news/10-ways-technology-is-changing-healthcare/).
In addition, the level of care that was received previously may have been very different than what is practiced today 
(https://prod.mgma.com/articles/from-error-to-excellence-redesigning-care-whole-system-quality). \\
Experience with a personal orthopaedic surgery has a minor overlap with having a knee surgery due to their similar nature.
In that sense, an orthopaedic patient may be more empathetic to those that have had surgeries recently. They may also impart
their first-hand understanding of such an experience and provide richer, more nuanced insights into the potential 
benefit and use of the application. 
\subsection{\textbf{Participant device literacy}}
In order to properly operate the application, there is a minimum level of device literacy needed. These questions specifically
target smartphone users. Some populations may not have equal access to the level of care (socioeconomic reasons) (https://www.cmaj.ca/content/185/6/E263) that is assumed 
for this application's use, and that is somewhat portrayed through the use of these questions. The purpose of these questions asserts the
need for further exploration if/when the target audience is widened and should further development of the application occurs. \\  %check with everyone if this sounds like bs

\subsection{\textbf{Participant perspective on privacy}}
This section outlines some private aspects involved in using the app. These statements 
should be confirmed with participants to ensure that a well-informed decision to participate and use it has been made.

\section{Session questionnaire}
\subsection{Create Account}
The account creation functionality is a similar one to that of existing apps. It uses a conventional flow that should be familiar to users.
The ease of this process is evaluated. 
Once a participant has created an account, they are directed to the Homescreen. This page has different capabilities based
on the type of user. \\
\subsection{Homescreen}
On a patient's screen, it displays a history of completed exercises and their dates. %TODO address this in discussion using input given by Kevin
On a PT's screen, it displays the names of their patients. Once they have selected a patient, the screen is populated
with the patient's activity. Both the patient's and PT's homescreen only serve as a record, displaying patient activity history. 
Since both patient and PTs pages look very similar, the main questions posed for this section are
UI perspectives. Drawing on both perspectives, these questions serve to validate the homescreen's minimal aesthetic,
confirming that it contains only essential information and is not visually cluttered. The same logic follows the progress
tab. \\

\subsection{Exercise}
Despite the UI being almost the same, there is a functional difference when the two types of users 
interact with the Exercises page which is addressed by the questions.\\
In the case of the patient, the questions aim to determine whether the elements are presented clearly and in a straightforward manner.
The example video under the complete exercise instructions was implemented at a later stage, necessitating investigation into whether
it was a helpful decision. \\
In the PT's case, an additional functionality is introduced. While the exercise list maintains the same appearance, clicking on the 
exercise reveals a text box which allows modifications to be made to the instructions. The intuitiveness and ease needed to be ascertained. \\

\subsection{Record}
PTs are unable to access this function and the page remains inaccessible. This functionality is not intended for them and was not
developed further. 
However, the patient's view provides additional features. The UI elements were assessed to confirm visual appropriateness. 
The following questions in the section evaluate whether the app successfully carries out its primary purpose. 


%TODO Feedback page 




%TODO post 


\end{document}




